\scanpage{4}
\deleted{par suite le deuxième membre de (8.1.4) est canoniquement dual à $\mathrm{Tors}^{n-1}(Y)\otimes_W k$. De même le premier membre est canoniquement \deleted{isomorphe à} \inserted{\handwritten{dual de}} $\mathrm{Tors}^{n+1}(X)\otimes_W k$, et enfin l'homomorphisme $u$ est transposé de l'homomorphisme}
\[
\deleted{u_{\handwritten{2}}'\mathbin{:}\mathrm{Tors}^{n-1}(Y)\otimes_W k\handwritten{\longrightarrow}\mathrm{Tors}^{n+1}(X)\otimes_W k}
\tag*{\deleted{(8.2.1)}}
\]
\deleted{induit par l'homomorphisme de Gysin}
\[
\deleted{H^{n-1}(Y)\handwritten{\longrightarrow}H^{n+1}(X)\;.}
\]

\deleted{Par suite, les termes $\mathrm{Ker}\,u_2$ et $\mathrm{Coker}\,u_2$ dans (8.1.3) peuvent s'interpréter \deleted{comme} respectivement comme \deleted{\unsure{Coker}} les duals de $\mathrm{Coker}\,u'_2$ et $\mathrm{Ker}\,u'_2$}

8.3. Considérons l'homomorphisme induit par l'inclusion \handwritten{$Y\longrightarrow X$,}
\[
u_{\deleted{\unsure{1}}0}\mathbin{:}H^{n-1}(X)\handwritten{\longrightarrow}H^{n-1}(Y)\;.
\]

\deleted{Sauf erreur,} \deleted{d}\handwritten{D}ans la situation analogue \deleted{sur le corps des complexes,} \inserted{\handwritten{en cohomologie $l$-adique ($l\ne p$)}} \uline{les \deleted{assertions} \inserted{\handwritten{théorème}} de Lefschetz} \deleted{dans son livre ``L'Analysis Situs ...''} \inserted{\handwritten{faible (pour les coefficients $Z/l^{\deleted{\unsure{m|n}}}Z$)}} impliquent que \uline{le conoyau} $E(Y)$ \uline{est sans torsion}, ce qui implique en particulier que $u_0$ induit un \uline{isomorphisme}
\[
\deleted{\unsure{H}}\handwritten{T}\mathrm{ors}^{n-1}(X)\handwritten{\overset{\sim}{\longrightarrow}}\mathrm{Tors}^{n-1}(Y)\;.
\tag*{(8.3.1)}
\]

Nous admettrons que ces relations restent valables dans le cas de la cohomologie cristalline. Alors l'homomorphisme (8.2.1) s'interprète comme l'homomorphisme
\[
u''_2\mathbin{:}\mathrm{Tors}^{n-1}(X)\otimes_W k\handwritten{\longrightarrow}\mathrm{Tors}^{n+1}(X)\otimes_W k
\tag*{(8.3.2)}
\]
induit par l'homomorphisme $H^{n-1}(X)\handwritten{\longrightarrow}H^{n+1}(X)$, cup-produit par la classe de polarisation de $X\handwritten{\subset}\mathrm P^r_k$. La chose importante pour nous est que ce morphisme (donc son Ker ) ne dépend plus du choix de $Y$, donc \deleted{aussi du} \deleted{\handwritten{\unsure{il en est de même}}} \inserted{\handwritten{le}} dernier terme de (8.1.3) \deleted{ne dépend} (qui est le dual de ce Ker) ne dépend plus de $Y$. Un autre point important est que, $E(Y)$ étant sans torsion, il a même rang que $E(Y)\otimes_W k$ intervenant dans (8.1.3) (et en particulier, si on a un endomorphisme de $E(Y)$, par exemple un frobenius, les valeurs propres de l'endomorphisme induit de $E(Y)\otimes k$ sont bien les réductions
% archive sheet number: 3
